% TOP_PROBLEM: {"problemId": "problem.montgomery-s-large-value-conjecture-for-dirichlet-polynomials", "canonicalTitle": "Montgomery's Large Value Conjecture for Dirichlet Polynomials", "releaseRank": 295, "primaryDomain": "number_theory_arithmetic_geometry", "primaryDomainLabel": "Number theory & arithmetic geometry", "displayStatus": "open_with_solved_subcases", "releaseStatus": "open_with_solved_subcases"}
% !TeX program = lualatex
% ATTEMPT_STATUS: unresolved
\documentclass[11pt]{article}
\usepackage[margin=25mm]{geometry}
\usepackage{fontspec}
\setmainfont{Latin Modern Roman}
\usepackage{amsmath,amssymb,mathtools}
\usepackage{unicode-math}
\setmathfont{Latin Modern Math}
\usepackage{xurl,hyperref}
\hypersetup{colorlinks=true,urlcolor=blue}
\setcounter{secnumdepth}{0}
\setlength{\parindent}{0pt}
\setlength{\parskip}{5pt}
\setlength{\emergencystretch}{3em}
\begin{document}
\section*{\texorpdfstring{Montgomery's Large Value Conjecture for Dirichlet Polynomials}{Montgomery's Large Value Conjecture for Dirichlet Polynomials}}\label{295}
\subsection{Definitions and mathematical statement}
Let $D(t)=\sum_{N<n\le2N}b_n n^{it}$ with $|b_n|\le1$, and let $W\subset[0,T]$ be one-separated: distinct elements differ by at least one. For fixed $\sigma>1/2$, assume $|D(t)|>N^\sigma$ on $W$. The target is
\[
 |W|\le T^{o(1)}N^{2-2\sigma},\qquad N\le T\le N^{O(1)}.
\]
Explicitly, with a fixed exponent range $T\le N^A$, the subpower notation asks for every $\eta>0$ a bound
$|W|\le C_{\sigma,A,\eta}T^\eta N^{2-2\sigma}$ uniformly in the coefficients and the separated set. An exponent range growing with $N$ is not the fixed-polynomial-range convention. The statement counts large values, not only an integral mean value.
\par\smallskip\textit{Formulation and concept references:} \href{https://arxiv.org/html/2405.20552}{[S1]}.

\subsection{Short English statement}
Can a bounded-coefficient Dirichlet polynomial take very large values at only the conjectured number of well-separated times in every polynomial-length interval?
\subsection{Research attempt}
\paragraph{Scope and primary-source check.}
This attempt was prepared by GPT-6 Astra Ultra on 27 September 2026.
Conjecture 1.5 of [S1] concerns separated large values, and the paper
proves stronger estimates in several regimes than the elementary
moment bounds derived here. We do not describe those elementary bounds
as current records. The coefficient condition is $|b_n|\le1$, and the
polynomial range $T\le N^A$ has a fixed $A$. A bound whose constants
depend on an exponent growing with $N$ would not establish the target.

Take $N$ to be a positive integer, as in the standard formulation.
If $\sigma\ge1$, the strict inequality $|D(t)|>N^\sigma$ is impossible,
because $|D(t)|\le N$. The substantive range is $1/2<\sigma<1$.
We prove the conjectured order, up to subpowers, when $T$ is comparable
to $N$, and give a family of higher-moment estimates for longer intervals.
The transition from integral means to separated samples is proved
explicitly instead of assuming that point values have positive measure.

\paragraph{An elementary mean-square estimate.}
Let $F(t)=\sum_{m\le M}a_m m^{it}$, with arbitrary complex coefficients,
and let $I$ be an interval of length $L$. Expanding and integrating,
the diagonal contributes $L\sum|a_m|^2$. For $m\ne l$,
\[
 \left|\int_I(m/l)^{it}\,dt\right|
                  \le\frac2{|\log(m/l)|}
                  \le\frac{2M}{|m-l|}.                 \tag{1}
\]
The last inequality follows by integrating $1/x$ between $m$ and $l$.
For each positive difference $h$, Cauchy--Schwarz gives
$\sum_m|a_ma_{m+h}|\le\sum_m|a_m|^2$, extending coefficients by zero.
Summing $1/h$ and both orders of the pair yields
\[
           \int_I|F(t)|^2\,dt
                   \le\bigl(L+C M\log(2M)\bigr)\sum_m|a_m|^2.
                                                               \tag{2}
\]
The bound is uniform in the location of the interval: translating it
only multiplies coefficients by unimodular factors. The logarithmic
loss is harmless for subpower purposes but is not suppressed before
it has been accounted for.

\paragraph{A separated-sample inequality.}
For a continuously differentiable complex function $f$ on
$J=[t_0-1/2,t_0+1/2]$, the fundamental theorem of calculus applied
to $|f|^2$ gives, for each $s\in J$,
\[
            |f(t_0)|^2\le|f(s)|^2+2\int_J|f f'|.
\]
Integrating over $s$ proves
\[
               |f(t_0)|^2\le\int_J|f|^2+2\int_J|f f'|. \tag{3}
\]
If $W$ is one-separated, the interiors of these length-one intervals
are disjoint. Summing (3) and using Cauchy--Schwarz gives
\[
 \sum_{t\in W}|f(t)|^2
 \le\int_{-1/2}^{T+1/2}|f|^2
 +2\left(\int_{-1/2}^{T+1/2}|f|^2
                \int_{-1/2}^{T+1/2}|f'|^2\right)^{1/2}. \tag{4}
\]
The slight enlargement of the interval handles samples at zero and
$T$. Separation is used here, and cannot be dropped.

Apply this to $f(t)=N^{-it}D(t)$, which has the same modulus as $D$.
Its derivative has coefficients $i b_n\log(n/N)$, bounded in modulus
by $\log2$. Estimate (2) applies equally to $f$ and $f'$, since
the common phase $N^{-it}$ cancels in the expanded integral.
We obtain
\[
         \sum_{t\in W}|D(t)|^2\le C(T+N\log(2N))N,
\]
and hence
\[
              |W|\le C(T+N\log(2N))N^{1-2\sigma}.       \tag{5}
\]
In particular, for $N\le T\le C_0N$, this is
$O_{C_0}(N^{2-2\sigma}\log(2N))$, which implies the requested
$O_{\sigma,C_0,\eta}(T^\eta N^{2-2\sigma})$ for every $\eta>0$.
This proves a genuine interval-range special case.

\paragraph{Integer moments and their coefficient growth.}
For a fixed positive integer $k$, expand
\[
 D(t)^k=\sum_m a_m m^{it},\qquad
 a_m=\sum_{\substack{N<n_1,\ldots,n_k\le2N\\n_1\cdots n_k=m}}
                                      b_{n_1}\cdots b_{n_k}.
\]
The support lies in $N^k<m\le(2N)^k$.
Let $r_m$ count the ordered tuples in this sum. Then
$|a_m|\le r_m\le d_k(m)$ and $\sum_mr_m=N^k$, so
\[
                 \sum_m|a_m|^2\le\max_{m\le(2N)^k}d_k(m)\,N^k.
                                                               \tag{6}
\]
For fixed $k$ and every $\epsilon>0$, the elementary divisor estimate
$d_k(m)\le C_{k,\epsilon}m^\epsilon$ follows from
\[
                 d_k(m)=\prod_{p^a\parallel m}\binom{a+k-1}{k-1}.
\]
For sufficiently large primes, $\binom{a+k-1}{k-1}\le k^a\le p^{\epsilon a}$.
For the finitely many smaller primes, the polynomial in $a$ is at most
a prime-dependent constant times $p^{\epsilon a}$; their constants
have a finite product. Renaming $\epsilon$ to absorb the fixed factor
$k$ in $m\le(2N)^k$, (6) gives
\[
                      \sum_m|a_m|^2\le C_{k,\epsilon}N^{k+\epsilon}.
                                                               \tag{7}
\]
This is where fixed $k$ matters. Constants in this proof are not
uniform when $k$ is allowed to grow with $N$.

Apply (4) to $f(t)=N^{-ikt}D(t)^k$.
Its derivative has multiplier $\log(m/N^k)$, bounded by $k\log2$.
Equations (2), (4), and (7) yield, after absorbing logarithms,
\[
        |W|\le C_{k,\epsilon}
                 (T+N^k)N^{k-2k\sigma+\epsilon}.        \tag{8}
\]
All coefficients and all separated sets allowed by the problem are
covered. The argument has not assumed random signs or an average over
the coefficients.

\paragraph{The exponent loss remains polynomial.}
Let $T=N^A$ with fixed $A\ge1$. Formula (8) has exponent
\[
                       \max\{A,k\}+k(1-2\sigma)+\epsilon.\tag{9}
\]
At $A=1$, $k=1$ gives the conjectured $2-2\sigma$.
At $A=2$, $\sigma=3/4$, even the best integer choices in this family
give exponent one, attained at $k=2$; the target exponent is $1/2$.
For $k=1$ the exponent is $3/2$, and for $k\ge2$ it is $k/2$.
Thus merely taking higher integer moments in this manner does not
close that representative gap. An $N^{1/2}$ loss cannot be renamed
$T^{o(1)}$: for $T=N^2$, it is $T^{1/4}$, which fails the requested
bound for sufficiently small fixed $\eta$.

This calculation diagnoses this particular moment argument, not the
limits of the stronger methods in [S1]. It also shows why invoking an
unspecified ``high moment'' without tracking its constants and length
can conceal the main difficulty.

\paragraph{A phase-coherence countertest.}
Choose $b_n=n^{-it_0}$. Then $D(t_0)=N$ exactly.
For $|s|\le1/2$, after removing the common phase $N^{is}$,
\[
 \operatorname{Re}\sum_{N<n\le2N}(n/N)^{is}
                         \ge N\cos((\log2)/2).
\]
Thus a high peak persists on an interval of fixed width.
The coefficients are fully admissible, so a proof cannot assume
pointwise cancellation at each time or treat all phases as independent.
This example gives only a bounded number of one-separated samples
in that interval and does not contradict the conjectured count.
Constructing many mutually compatible peaks is a separate issue.

\paragraph{Verification and final status.}
The local checks expand small powers exactly, count the product
multiplicities in (6), and verify the algebra in (9). Numerical phase
samples are labeled only as sanity checks of the explicit example.
The proof of (4) supplies the otherwise missing link between integral
and pointwise large-value statements.

\textbf{UNRESOLVED.} The result proved here is the requested estimate
for $T\asymp N$, together with the uniform fixed-moment family (8).
The first missing step is a separated-sample estimate eliminating the
extra polynomial dependence on $T/N$ in the longer ranges. Concrete
targets are correlations among large-value times, cancellation in the
off-diagonal moment sum before absolute values, and estimates for product
coefficients sharper than their maximum multiplicity. The arbitrary
fixed-polynomial-range conjecture is not established.

\subsection{Sources}
\begin{itemize}
\item[S1]Larry Guth and James Maynard, New large value estimates for Dirichlet polynomials, Annals of Mathematics 203 (2026), 623-675, Conjecture 1.5.
\url{https://arxiv.org/html/2405.20552}.
\textit{Formulation source.}
\end{itemize}
\end{document}
